% GENERATED FILE. Edit canonical lessons, then run npm run generate:books.
% canonical-source-hash: fnv1a64:0e543c014cc9dda1
% canonical-lessons: SA-C110-girih, SA-C110-prapatah, SA-C110-tarangah, SA-C110-vidyut, SA-C110-garjanam

\chapter{Mountain, Waterfall, Wave, Lightning, Thunder}
\label{ch:sa-mountain-waterfall-wave-lightning-thunder}
\hlchaptermodality{110}

\begin{chapteropening}
\textbf{By the end of this chapter:} \emph{I can say a mountain, a waterfall, a wave, lightning and thunder.}

Complete the last lesson of chapter 110: i can say a mountain, a waterfall, a wave, lightning and thunder.
\end{chapteropening}

\section[giriḥ]{\sk{गिरिः} (giriḥ) — a mountain, a hill}

\label{lesson:SA-C110-girih}



\begin{quote}
Before the new one: say the Sanskrit for a lemon, then the Sanskrit for a cake.
\end{quote}

\subsection*{You'll want to know: \sk{गिरिः}}
\textbf{\sk{गिरिः}} — \emph{giriḥ} — \textquotedblleft{}a mountain, a hill\textquotedblright{}.

\begin{grammarlens}[title={What you've built}]
One more word for this chapter.
\end{grammarlens}

\subsection*{Your turn}
\begin{itemize}
\raggedright
  \item \emph{Say it:} \emph{giriḥ}
  \item \emph{Say it:} \emph{giriḥ}, once more
  \item \emph{Say it:} say \emph{giriḥ}
  \item \emph{Recall:} say the Sanskrit for a lemon, then the Sanskrit for a cake, then say \emph{giriḥ} again
\end{itemize}

\begin{tcolorbox}[breakable,colback=teal!4,colframe=teal!35!black,title={Before you move on}]
What does \sk{गिरिः} mean? (\textquotedblleft{}A mountain, a hill\textquotedblright{}.) Say it once more.
\end{tcolorbox}
\section[prapātaḥ]{\sk{प्रपातः} (prapātaḥ) — a waterfall}

\label{lesson:SA-C110-prapatah}



\begin{quote}
Before the new one: say the Sanskrit for a cake, then the Sanskrit for a mountain.
\end{quote}

\subsection*{You'll want to know: \sk{प्रपातः}}
\textbf{\sk{प्रपातः}} — \emph{prapātaḥ} — \textquotedblleft{}a waterfall\textquotedblright{}.

\begin{grammarlens}[title={What you've built}]
One more word for this chapter.
\end{grammarlens}

\subsection*{Your turn}
\begin{itemize}
\raggedright
  \item \emph{Say it:} \emph{prapātaḥ}
  \item \emph{Say it:} \emph{prapātaḥ}, once more
  \item \emph{Say it:} say \emph{prapātaḥ}
  \item \emph{Recall:} say the Sanskrit for a cake, then the Sanskrit for a mountain, then say \emph{prapātaḥ} again
\end{itemize}

\begin{tcolorbox}[breakable,colback=teal!4,colframe=teal!35!black,title={Before you move on}]
What does \sk{प्रपातः} mean? (\textquotedblleft{}A waterfall\textquotedblright{}.) Say it once more.
\end{tcolorbox}
\section[taraṅgaḥ]{\sk{तरंगः} (taraṅgaḥ) — a wave}

\label{lesson:SA-C110-tarangah}



\begin{quote}
Before the new one: say the Sanskrit for a mountain, then the Sanskrit for a waterfall.
\end{quote}

\subsection*{You'll want to know: \sk{तरंगः}}
\textbf{\sk{तरंगः}} — \emph{taraṅgaḥ} — \textquotedblleft{}a wave\textquotedblright{}.

\begin{grammarlens}[title={What you've built}]
One more word for this chapter.
\end{grammarlens}

\subsection*{Your turn}
\begin{itemize}
\raggedright
  \item \emph{Say it:} \emph{taraṅgaḥ}
  \item \emph{Say it:} \emph{taraṅgaḥ}, once more
  \item \emph{Say it:} say \emph{taraṅgaḥ}
  \item \emph{Recall:} say the Sanskrit for a mountain, then the Sanskrit for a waterfall, then say \emph{taraṅgaḥ} again
\end{itemize}

\begin{tcolorbox}[breakable,colback=teal!4,colframe=teal!35!black,title={Before you move on}]
What does \sk{तरंगः} mean? (\textquotedblleft{}A wave\textquotedblright{}.) Say it once more.
\end{tcolorbox}
\section[vidyut]{\sk{विद्युत्} (vidyut) — lightning}

\label{lesson:SA-C110-vidyut}



\begin{quote}
Before the new one: say the Sanskrit for a waterfall, then the Sanskrit for a wave.
\end{quote}

\subsection*{You'll want to know: \sk{विद्युत्}}
\textbf{\sk{विद्युत्}} — \emph{vidyut} — \textquotedblleft{}lightning\textquotedblright{}.

\begin{grammarlens}[title={What you've built}]
One more word for this chapter.
\end{grammarlens}

\subsection*{Your turn}
\begin{itemize}
\raggedright
  \item \emph{Say it:} \emph{vidyut}
  \item \emph{Say it:} \emph{vidyut}, once more
  \item \emph{Say it:} say \emph{vidyut}
  \item \emph{Recall:} say the Sanskrit for a waterfall, then the Sanskrit for a wave, then say \emph{vidyut} again
\end{itemize}

\begin{tcolorbox}[breakable,colback=teal!4,colframe=teal!35!black,title={Before you move on}]
What does \sk{विद्युत्} mean? (\textquotedblleft{}Lightning\textquotedblright{}.) Say it once more.
\end{tcolorbox}
\section[garjanam]{\sk{गर्जनम्} (garjanam) — thunder}

\label{lesson:SA-C110-garjanam}



\begin{quote}
Before the new one: say the Sanskrit for a wave, then the Sanskrit for lightning.
\end{quote}

\subsection*{You'll want to know: \sk{गर्जनम्}}
\textbf{\sk{गर्जनम्}} — \emph{garjanam} — \textquotedblleft{}thunder\textquotedblright{}.

\begin{grammarlens}[title={What you've built}]
One more word for this chapter.
\end{grammarlens}

\subsection*{Your turn}
\begin{itemize}
\raggedright
  \item \emph{Say it:} \emph{garjanam}
  \item \emph{Say it:} \emph{garjanam}, once more
  \item \emph{Say it:} say \emph{garjanam}
  \item \emph{Recall:} say the Sanskrit for a wave, then the Sanskrit for lightning, then say \emph{garjanam} again
\end{itemize}

\begin{tcolorbox}[breakable,colback=teal!4,colframe=teal!35!black,title={Before you move on}]
What does \sk{गर्जनम्} mean? (\textquotedblleft{}Thunder\textquotedblright{}.) Say it once more.
\end{tcolorbox}
